
\section{Der Funktionen-Raum}
\label{sec:1.1}

\begin{definition}{Riemann-Integrierbarkeit}
    Eine Funktion \(f : [a,b] \to \mathbb{C}\) heißt \emph{Riemann-integrierbar} auf \([a,b]\),
    falls \(\Re(f)\) und \(\Im(f)\) Riemann-integrierbar sind.

    Man definiert das Integral komplexwertiger Funktionen durch
    \[
    \int_a^b f(x)\,dx
    := \int_a^b \Re f(x)\,dx \;+\; i \int_a^b \Im f(x)\,dx.
    \]
\end{definition}

\begin{bemerkungbox}
    \begin{enumerate}
        \item Uneigentliche Riemann-Integrale lassen sich analog für komplexwertige Funktionen definieren.
        \item Die Rechenregeln des reellen Riemann-Integrals übertragen sich auf komplexwertige Integrale.
    \end{enumerate}
\end{bemerkungbox}

\begin{definition}{Stückweise-Stetigkeit}
    Eine Funktion \(f : [a,b] \to K\) (mit \(K = \mathbb{R}\) oder \(\mathbb{C}\)) heißt
    \emph{stückweise stetig}, falls:

    \begin{enumerate}
        \item \(f\) auf \([a,b]\) beschränkt ist und nur endlich viele Unstetigkeitsstellen besitzt,
        \item an jeder Unstetigkeitsstelle \(\xi\) die Grenzwerte
        \[
            f(\xi^\pm) := \lim_{h\to 0^\pm} f(\xi + h)
        \]
        existieren.
    \end{enumerate}

    Für \(\xi \in (a,b)\) setzt man zusätzlich die Konvention

    \[
    f(\xi) := \frac{f(\xi^-) + f(\xi^+)}{2}.
    \]
\end{definition}

\begin{bemerkungbox}
\begin{itemize}
    \item Stückweise stetige Funktionen sind Riemann-integrierbar.
    \item Die Menge aller stückweise stetigen Funktionen bildet den Vektorraum \(R[a,b]\).
\end{itemize}
\end{bemerkungbox}

\begin{definition}{Sesquilinearform}
    Für \(f,g \in R[a,b]\) definiert man
    \[
        (f,g) := \int_a^b f(x)\,\overline{g(x)}\,dx.
    \]
    Diese Abbildung ist wohldefiniert, da \(f,g \in R[a,b]\) und somit auch
    \(f\cdot g \in R[a,b]\).
\end{definition}

\begin{bemerkungbox}
    Für alle \(f,f_1,f_2,g,g_1,g_2 \in R[a,b]\) und Skalare \(\alpha,\beta\) gilt:
    \begin{enumerate}
        \item Linearität im ersten Argument:
        \[
        (\alpha f_1 + \beta f_2, g)
        = \alpha (f_1,g) + \beta (f_2,g).
        \]
        \item Linearität im zweiten Argument:
        \[
        (f, \alpha g_1 + \beta g_2)
        = \alpha (f,g_1) + \beta (f,g_2).
        \]
        \item Symmetrieeigenschaft:
        \[
        (f,g) = \overline{(g,f)}.
        \]
        \item Semidefinitheit:
        \[
        (f,f) = \int_a^b |f(x)|^2\,dx \ge 0.
        \]
    \end{enumerate}
\end{bemerkungbox}

\begin{definition}{Skalarprodukt}
    Sei \(V\) ein Vektorraum über \(K\) (mit \(K = \mathbb{R}\) oder \(\mathbb{C}\)).
    Eine Abbildung \((\cdot,\cdot): V \times V \to K\) heißt \emph{Skalarprodukt}, wenn:
    \begin{enumerate}
        \item Symmetrie:
        \[
            (v,u) = \overline{(u,v)}.
        \]
        \item Linearität:
        \[
            (\alpha v, u) = \alpha (v,u), \qquad (v, u+w) = (v,u) + (v,w).
        \]
        \item Positivdefinitheit:
        \[
            (v,v) \ge 0, \quad (v,v)=0 \iff v=0.
        \]
    \end{enumerate}
\end{definition}

\begin{lemma}
    Für alle \(f,g \in R[a,b]\) gilt die Cauchy–Schwarz-Ungleichung:
    \[
        |(f,g)|^2 \le (f,f)\,(g,g).
    \]
\end{lemma}

\begin{definition}{L\textsuperscript{2}-Norm}
    Die durch das Skalarprodukt induzierte Norm auf \(R[a,b]\) ist definiert durch
    \[
        \|f\|_2 := \left( \int_a^b |f(x)|^2\,dx \right)^{1/2}.
    \]
\end{definition}

\begin{bemerkungbox}
    Die L\textsuperscript{2}-Norm erfüllt:
    \begin{enumerate}
        \item Definitheit: \(\|f\|_2 = 0 \iff f = 0\).
        \item Homogenität: \(\|\alpha f\|_2 = |\alpha|\,\|f\|_2\).
        \item Dreiecksungleichung:
        \[
            \|f+g\|_2 \le \|f\|_2 + \|g\|_2.
        \]
    \end{enumerate}
\end{bemerkungbox}


\begin{definition}{L\textsuperscript{2}-Konvergenz}
    Eine Folge \((f_n)\subset R[a,b]\) konvergiert \emph{im quadratischen Mittel}
    gegen \(f\in R[a,b]\), wenn
    \[
        \|f_n - f\|_2 \;\longrightarrow\; 0.
    \]

    Äquivalent:
    \[
        \int_a^b |f_n(x) - f(x)|^2\,dx \;\to\; 0.
    \]
\end{definition}

\begin{bemerkungbox}
    Für jede Folge \(f_n\) gilt:
    \[
        \|f_n - f\|_\infty^2 \le (b-a)\,\|f_n - f\|_2^2.
    \]

    Daraus folgt:
    \[
        \|f_n - f\|_\infty \to 0 \;\Rightarrow\; \|f_n - f\|_2 \to 0,
    \]
    aber die Umkehrung gilt im Allgemeinen nicht.
\end{bemerkungbox}

\begin{beispielbox}
    Betrachte \(f_n(x) = x^n\) auf \([-1,1]\). Dann gilt:
    \[
        \|f_n\|_2^2 = 2 \int_0^1 x^{2n}\,dx = \frac{2}{2n+1} \longrightarrow 0,
    \]
    also \(f_n \to 0\) in L\(^2\).

    Punktweise gilt jedoch:
    \[
        f_n(1) = 1 \quad\text{für alle } n,
    \]
    also keine punktweise Konvergenz auf ganz \([-1,1]\).
\end{beispielbox}

\begin{definition}{L\textsuperscript{p}-Norm}
    Für \(p \in [1,\infty)\) und \(f : [a,b] \to \mathbb{R}\) definiert man:
    \[
        \|f\|_p := \left( \int_a^b |f(x)|^p\,dx \right)^{1/p}.
    \]
    Für \(p=2\) erhält man die L\(^2\)-Norm.
\end{definition}

\begin{bemerkungbox}
    Der Raum \(R[a,b]\) ist bezüglich der L\(^2\)-Norm nicht vollständig.
    Es existieren Cauchy-Folgen in \(R[a,b]\), die keinen Grenzwert in
    \(R[a,b]\) besitzen. Die Vervollständigung führt zum Raum \(L^2[a,b]\).
\end{bemerkungbox}

\begin{definition}{Orthogonalität}
    Für \(f,g \in R[a,b]\) heißt \(f\) orthogonal zu \(g\), wenn
    \[
        (f,g) = 0.
    \]
    Eine Menge \(S \subset R[a,b]\) heißt \emph{orthogonal}, wenn für alle
    \(f_i,f_j \in S\):
    \[
        (f_i,f_j) =
        \begin{cases}
            0, & i \ne j,\\
            \|f_i\|_2^2, & i = j.
        \end{cases}
    \]
\end{definition}


\begin{satz}{Orthogonalität trigonometrischer Funktionen}
    Die Funktionen
    \[
        c_0(x)=1,\quad c_k(x)=\cos(kx),\quad s_k(x)=\sin(kx),\quad k\in\mathbb{N},
    \]
    bilden auf \(R[0,2\pi]\) ein Orthogonalsystem bezüglich des L\(^2\)-Skalarprodukts.
    Es gilt:
    \[
        \int_0^{2\pi} c_k(x)c_l(x)\,dx = \pi\,\delta_{kl},\qquad
        \int_0^{2\pi} s_k(x)s_l(x)\,dx = \pi\,\delta_{kl},
    \]
    \[
        \int_0^{2\pi} c_k(x)s_l(x)\,dx = 0.
    \]
\end{satz}

\begin{definition}{Periodische Funktion}
    Eine Funktion \(f:\mathbb{R}\to K\) heißt \(L\)-periodisch, wenn
    \[
        f(x+L)=f(x)\quad\forall x\in\mathbb{R}.
    \]
\end{definition}
